% Options for packages loaded elsewhere
\PassOptionsToPackage{unicode}{hyperref}
\PassOptionsToPackage{hyphens}{url}
\documentclass[
]{article}
\usepackage{xcolor}
\usepackage[margin=1in]{geometry}
\usepackage{amsmath,amssymb}
\setcounter{secnumdepth}{-\maxdimen} % remove section numbering
\usepackage{iftex}
\ifPDFTeX
  \usepackage[T1]{fontenc}
  \usepackage[utf8]{inputenc}
  \usepackage{textcomp} % provide euro and other symbols
\else % if luatex or xetex
  \usepackage{unicode-math} % this also loads fontspec
  \defaultfontfeatures{Scale=MatchLowercase}
  \defaultfontfeatures[\rmfamily]{Ligatures=TeX,Scale=1}
\fi
\usepackage{lmodern}
\ifPDFTeX\else
  % xetex/luatex font selection
\fi
% Use upquote if available, for straight quotes in verbatim environments
\IfFileExists{upquote.sty}{\usepackage{upquote}}{}
\IfFileExists{microtype.sty}{% use microtype if available
  \usepackage[]{microtype}
  \UseMicrotypeSet[protrusion]{basicmath} % disable protrusion for tt fonts
}{}
\makeatletter
\@ifundefined{KOMAClassName}{% if non-KOMA class
  \IfFileExists{parskip.sty}{%
    \usepackage{parskip}
  }{% else
    \setlength{\parindent}{0pt}
    \setlength{\parskip}{6pt plus 2pt minus 1pt}}
}{% if KOMA class
  \KOMAoptions{parskip=half}}
\makeatother
\usepackage{color}
\usepackage{fancyvrb}
\newcommand{\VerbBar}{|}
\newcommand{\VERB}{\Verb[commandchars=\\\{\}]}
\DefineVerbatimEnvironment{Highlighting}{Verbatim}{commandchars=\\\{\}}
% Add ',fontsize=\small' for more characters per line
\usepackage{framed}
\definecolor{shadecolor}{RGB}{248,248,248}
\newenvironment{Shaded}{\begin{snugshade}}{\end{snugshade}}
\newcommand{\AlertTok}[1]{\textcolor[rgb]{0.94,0.16,0.16}{#1}}
\newcommand{\AnnotationTok}[1]{\textcolor[rgb]{0.56,0.35,0.01}{\textbf{\textit{#1}}}}
\newcommand{\AttributeTok}[1]{\textcolor[rgb]{0.13,0.29,0.53}{#1}}
\newcommand{\BaseNTok}[1]{\textcolor[rgb]{0.00,0.00,0.81}{#1}}
\newcommand{\BuiltInTok}[1]{#1}
\newcommand{\CharTok}[1]{\textcolor[rgb]{0.31,0.60,0.02}{#1}}
\newcommand{\CommentTok}[1]{\textcolor[rgb]{0.56,0.35,0.01}{\textit{#1}}}
\newcommand{\CommentVarTok}[1]{\textcolor[rgb]{0.56,0.35,0.01}{\textbf{\textit{#1}}}}
\newcommand{\ConstantTok}[1]{\textcolor[rgb]{0.56,0.35,0.01}{#1}}
\newcommand{\ControlFlowTok}[1]{\textcolor[rgb]{0.13,0.29,0.53}{\textbf{#1}}}
\newcommand{\DataTypeTok}[1]{\textcolor[rgb]{0.13,0.29,0.53}{#1}}
\newcommand{\DecValTok}[1]{\textcolor[rgb]{0.00,0.00,0.81}{#1}}
\newcommand{\DocumentationTok}[1]{\textcolor[rgb]{0.56,0.35,0.01}{\textbf{\textit{#1}}}}
\newcommand{\ErrorTok}[1]{\textcolor[rgb]{0.64,0.00,0.00}{\textbf{#1}}}
\newcommand{\ExtensionTok}[1]{#1}
\newcommand{\FloatTok}[1]{\textcolor[rgb]{0.00,0.00,0.81}{#1}}
\newcommand{\FunctionTok}[1]{\textcolor[rgb]{0.13,0.29,0.53}{\textbf{#1}}}
\newcommand{\ImportTok}[1]{#1}
\newcommand{\InformationTok}[1]{\textcolor[rgb]{0.56,0.35,0.01}{\textbf{\textit{#1}}}}
\newcommand{\KeywordTok}[1]{\textcolor[rgb]{0.13,0.29,0.53}{\textbf{#1}}}
\newcommand{\NormalTok}[1]{#1}
\newcommand{\OperatorTok}[1]{\textcolor[rgb]{0.81,0.36,0.00}{\textbf{#1}}}
\newcommand{\OtherTok}[1]{\textcolor[rgb]{0.56,0.35,0.01}{#1}}
\newcommand{\PreprocessorTok}[1]{\textcolor[rgb]{0.56,0.35,0.01}{\textit{#1}}}
\newcommand{\RegionMarkerTok}[1]{#1}
\newcommand{\SpecialCharTok}[1]{\textcolor[rgb]{0.81,0.36,0.00}{\textbf{#1}}}
\newcommand{\SpecialStringTok}[1]{\textcolor[rgb]{0.31,0.60,0.02}{#1}}
\newcommand{\StringTok}[1]{\textcolor[rgb]{0.31,0.60,0.02}{#1}}
\newcommand{\VariableTok}[1]{\textcolor[rgb]{0.00,0.00,0.00}{#1}}
\newcommand{\VerbatimStringTok}[1]{\textcolor[rgb]{0.31,0.60,0.02}{#1}}
\newcommand{\WarningTok}[1]{\textcolor[rgb]{0.56,0.35,0.01}{\textbf{\textit{#1}}}}
\usepackage{longtable,booktabs,array}
\newcounter{none} % for unnumbered tables
\usepackage{calc} % for calculating minipage widths
% Correct order of tables after \paragraph or \subparagraph
\usepackage{etoolbox}
\makeatletter
\patchcmd\longtable{\par}{\if@noskipsec\mbox{}\fi\par}{}{}
\makeatother
% Allow footnotes in longtable head/foot
\IfFileExists{footnotehyper.sty}{\usepackage{footnotehyper}}{\usepackage{footnote}}
\makesavenoteenv{longtable}
\usepackage{graphicx}
\makeatletter
\newsavebox\pandoc@box
\newcommand*\pandocbounded[1]{% scales image to fit in text height/width
  \sbox\pandoc@box{#1}%
  \Gscale@div\@tempa{\textheight}{\dimexpr\ht\pandoc@box+\dp\pandoc@box\relax}%
  \Gscale@div\@tempb{\linewidth}{\wd\pandoc@box}%
  \ifdim\@tempb\p@<\@tempa\p@\let\@tempa\@tempb\fi% select the smaller of both
  \ifdim\@tempa\p@<\p@\scalebox{\@tempa}{\usebox\pandoc@box}%
  \else\usebox{\pandoc@box}%
  \fi%
}
% Set default figure placement to htbp
\def\fps@figure{htbp}
\makeatother
\setlength{\emergencystretch}{3em} % prevent overfull lines
\providecommand{\tightlist}{%
  \setlength{\itemsep}{0pt}\setlength{\parskip}{0pt}}
\usepackage{bookmark}
\IfFileExists{xurl.sty}{\usepackage{xurl}}{} % add URL line breaks if available
\urlstyle{same}
\hypersetup{
  pdftitle={Analyse Statistique des Données Physiologiques et Anthropométriques},
  pdfauthor={Ghada},
  hidelinks,
  pdfcreator={LaTeX via pandoc}}

\title{Analyse Statistique des Données Physiologiques et
Anthropométriques}
\author{Ghada}
\date{2026-05-14}

\begin{document}
\maketitle

{
\setcounter{tocdepth}{3}
\tableofcontents
}
\begin{center}\rule{0.5\linewidth}{0.5pt}\end{center}

\section{Introduction}\label{introduction}

\subsection{Contexte du projet}\label{contexte-du-projet}

Dans le domaine des sciences du sport, de la santé publique et de la
biométrie, l'analyse statistique des données physiologiques et
anthropométriques joue un rôle fondamental dans la compréhension des
déterminants de la condition physique humaine. Les indicateurs tels que
l'âge, l'indice de masse corporelle (IMC), le pourcentage de masse
grasse, la force musculaire, l'endurance et la capacité cardiovasculaire
sont largement utilisés pour évaluer l'état de santé et le niveau de
forme physique des individus.

Cependant, ces variables présentent souvent des distributions complexes,
influencées par des facteurs biologiques, sociaux et comportementaux,
notamment le sexe. L'étude statistique rigoureuse de ces données
nécessite donc l'utilisation conjointe de méthodes descriptives, de
tests de normalité et de tests paramétriques et non paramétriques
permettant d'identifier les relations significatives entre les
variables.

Dans ce contexte, le présent projet s'inscrit dans une démarche
d'analyse statistique approfondie d'un jeu de données réel contenant des
informations démographiques, anthropométriques et des mesures de
performance physique. Ce travail vise à appliquer différentes techniques
statistiques afin de valider des hypothèses, d'évaluer la structure des
données et de comparer les groupes selon des variables pertinentes.

Ce projet constitue également une application pratique des concepts
théoriques abordés dans le cadre du cours de statistique, en mettant
l'accent sur la rigueur méthodologique, l'interprétation correcte des
résultats et la justification du choix des tests statistiques utilisés.

\subsection{Objectifs du projet}\label{objectifs-du-projet}

L'objectif principal de ce projet est de réaliser une analyse
statistique complète et structurée d'un ensemble de données
physiologiques et anthropométriques, en vue de mieux comprendre les
relations entre les caractéristiques individuelles et les performances
physiques.

Plus précisément, ce projet vise à :

\begin{itemize}
\tightlist
\item
  \textbf{Explorer et décrire} les données à l'aide de statistiques
  descriptives afin d'identifier les tendances centrales, la dispersion
  et la structure globale des variables étudiées.
\item
  \textbf{Vérifier les hypothèses de normalité} des variables
  quantitatives à l'aide de tests statistiques appropriés, notamment le
  test de Shapiro--Wilk, afin de justifier le choix des méthodes
  paramétriques ou non paramétriques.
\item
  \textbf{Détecter et analyser} les valeurs aberrantes susceptibles
  d'influencer les résultats statistiques.
\item
  \textbf{Comparer les groupes} définis par la variable sexe à l'aide de
  tests paramétriques (tests t, ANOVA, tests de proportion) et non
  paramétriques (tests de Wilcoxon, tests des signes), selon la nature
  et la distribution des données.
\item
  \textbf{Étudier les relations} entre les variables à travers l'analyse
  des corrélations et l'évaluation de leur significativité statistique.
\end{itemize}

À travers ces objectifs, ce projet cherche à démontrer la capacité des
outils statistiques à fournir des résultats fiables et interprétables,
tout en mettant en évidence l'importance du choix méthodologique dans
l'analyse des données réelles.

\subsection{Description des données}\label{description-des-donnuxe9es}

Les données utilisées dans ce projet proviennent d'un ensemble de
mesures physiologiques, anthropométriques et de performance physique
collectées auprès d'un échantillon d'individus. Chaque observation
correspond à un individu et regroupe des informations démographiques,
des caractéristiques corporelles ainsi que des indicateurs de condition
physique.

Les variables étudiées ont été sélectionnées afin de couvrir différents
aspects de la condition physique, notamment la composition corporelle,
la force musculaire, l'endurance musculaire, la flexibilité et la
capacité cardiorespiratoire. La présence de variables qualitatives et
quantitatives permet de mener des analyses statistiques variées,
incluant des comparaisons de groupes et des tests d'hypothèses.

{\def\LTcaptype{none} % do not increment counter
\begin{longtable}[]{@{}
  >{\raggedright\arraybackslash}p{(\linewidth - 4\tabcolsep) * \real{0.3448}}
  >{\raggedright\arraybackslash}p{(\linewidth - 4\tabcolsep) * \real{0.2069}}
  >{\raggedright\arraybackslash}p{(\linewidth - 4\tabcolsep) * \real{0.4483}}@{}}
\toprule\noalign{}
\begin{minipage}[b]{\linewidth}\raggedright
Variable
\end{minipage} & \begin{minipage}[b]{\linewidth}\raggedright
Type
\end{minipage} & \begin{minipage}[b]{\linewidth}\raggedright
Description
\end{minipage} \\
\midrule\noalign{}
\endhead
\bottomrule\noalign{}
\endlastfoot
\texttt{sex} & Qualitative & Sexe de l'individu (M/F) \\
\texttt{age} & Quantitative & Âge de l'individu (années) \\
\texttt{bmi} & Quantitative & Indice de masse corporelle \\
\texttt{percent\_body\_fat} & Quantitative & Pourcentage de masse
grasse \\
\texttt{hand\_grip\_strength\_kg} & Quantitative & Force de préhension
(kg) \\
\texttt{sit\_ups\_count} & Quantitative & Nombre de sit-ups \\
\texttt{sit\_and\_reach\_cm} & Quantitative & Flexibilité (cm) \\
\texttt{vo2\_estimate\_ml\_per\_kg\_min} & Quantitative & VO₂ max estimé
(ml/kg/min) \\
\end{longtable}
}

\begin{center}\rule{0.5\linewidth}{0.5pt}\end{center}

\section{Importation et Préparation des
Données}\label{importation-et-pruxe9paration-des-donnuxe9es}

Cette section présente les étapes préliminaires nécessaires à l'analyse
statistique des données : chargement des bibliothèques, importation du
jeu de données, exploration initiale et étude de la distribution des
variables quantitatives. Ces étapes sont essentielles pour assurer la
qualité des analyses ultérieures et orienter le choix des méthodes
appropriées.

\subsection{Chargement des packages}\label{chargement-des-packages}

L'analyse des données a été réalisée à l'aide du logiciel R, en
utilisant plusieurs bibliothèques spécialisées :

\begin{itemize}
\tightlist
\item
  \texttt{openxlsx} et \texttt{readxl} : importation de fichiers Excel
\item
  \texttt{dplyr} et \texttt{tidyr} : manipulation et transformation des
  données
\item
  \texttt{ggplot2} et \texttt{gridExtra} : visualisation graphique
\item
  \texttt{psych} : statistiques descriptives
\item
  \texttt{corrplot} : visualisation des corrélations
\item
  \texttt{MASS} : méthodes statistiques avancées (régression pas à pas)
\item
  \texttt{naniar} : détection et visualisation des valeurs manquantes
\end{itemize}

\begin{Shaded}
\begin{Highlighting}[]
\FunctionTok{library}\NormalTok{(readxl)}
\FunctionTok{library}\NormalTok{(dplyr)}
\FunctionTok{library}\NormalTok{(ggplot2)}
\FunctionTok{library}\NormalTok{(corrplot)}
\end{Highlighting}
\end{Shaded}

\subsection{Importation et exploration
initiale}\label{importation-et-exploration-initiale}

Le jeu de données est importé à partir d'un fichier Excel à l'aide de la
fonction \texttt{read\_excel()}. Les fonctions \texttt{head()},
\texttt{summary()} et \texttt{str()} permettent d'obtenir un aperçu des
premières lignes, des statistiques descriptives générales et des
informations sur le type de chaque variable.

\begin{Shaded}
\begin{Highlighting}[]
\NormalTok{data }\OtherTok{\textless{}{-}} \FunctionTok{read\_excel}\NormalTok{(}\StringTok{"DATA.xlsx"}\NormalTok{)}

\FunctionTok{head}\NormalTok{(data)}
\end{Highlighting}
\end{Shaded}

\begin{verbatim}
## # A tibble: 6 x 10
##   participant_id sex     age measurement_date   bmi percent_body_fat
##            <dbl> <chr> <dbl> <chr>            <dbl>            <dbl>
## 1              1 M        46 2019-01-19        26.5             45  
## 2              2 M        25 2018-03-16        28.8             45  
## 3              3 M        53 2017-06-02        25.7             45  
## 4              4 M        46 2018-02-11        28               45  
## 5              5 F        32 2018-01-25        25               57.7
## 6              6 F        28 2018-11-06        24.7             60  
## # i 4 more variables: hand_grip_strength_kg <dbl>, sit_and_reach_cm <dbl>,
## #   sit_ups_count <dbl>, vo2_estimate_ml_per_kg_min <dbl>
\end{verbatim}

\begin{Shaded}
\begin{Highlighting}[]
\FunctionTok{dim}\NormalTok{(data)}
\end{Highlighting}
\end{Shaded}

\begin{verbatim}
## [1] 2000   10
\end{verbatim}

\begin{Shaded}
\begin{Highlighting}[]
\FunctionTok{str}\NormalTok{(data)}
\end{Highlighting}
\end{Shaded}

\begin{verbatim}
## tibble [2,000 x 10] (S3: tbl_df/tbl/data.frame)
##  $ participant_id            : num [1:2000] 1 2 3 4 5 6 7 8 9 10 ...
##  $ sex                       : chr [1:2000] "M" "M" "M" "M" ...
##  $ age                       : num [1:2000] 46 25 53 46 32 28 34 34 49 60 ...
##  $ measurement_date          : chr [1:2000] "2019-01-19" "2018-03-16" "2017-06-02" "2018-02-11" ...
##  $ bmi                       : num [1:2000] 26.5 28.8 25.7 28 25 24.7 24.7 27.8 25 24.6 ...
##  $ percent_body_fat          : num [1:2000] 45 45 45 45 57.7 60 60 45 45 45 ...
##  $ hand_grip_strength_kg     : num [1:2000] 47 28 37.7 48.9 20 21.2 28.9 38.2 41.9 34.1 ...
##  $ sit_and_reach_cm          : num [1:2000] 13.4 -1.2 4.5 9 14.8 13.1 3.9 17.1 9.9 14.6 ...
##  $ sit_ups_count             : num [1:2000] 13 19 18 20 23 27 17 12 12 15 ...
##  $ vo2_estimate_ml_per_kg_min: num [1:2000] 29.8 46.5 37.2 52.5 24.6 30.1 30 45.7 39.5 28.8 ...
\end{verbatim}

\begin{Shaded}
\begin{Highlighting}[]
\FunctionTok{colnames}\NormalTok{(data)}
\end{Highlighting}
\end{Shaded}

\begin{verbatim}
##  [1] "participant_id"             "sex"                       
##  [3] "age"                        "measurement_date"          
##  [5] "bmi"                        "percent_body_fat"          
##  [7] "hand_grip_strength_kg"      "sit_and_reach_cm"          
##  [9] "sit_ups_count"              "vo2_estimate_ml_per_kg_min"
\end{verbatim}

\begin{Shaded}
\begin{Highlighting}[]
\FunctionTok{sapply}\NormalTok{(data, class)}
\end{Highlighting}
\end{Shaded}

\begin{verbatim}
##             participant_id                        sex 
##                  "numeric"                "character" 
##                        age           measurement_date 
##                  "numeric"                "character" 
##                        bmi           percent_body_fat 
##                  "numeric"                  "numeric" 
##      hand_grip_strength_kg           sit_and_reach_cm 
##                  "numeric"                  "numeric" 
##              sit_ups_count vo2_estimate_ml_per_kg_min 
##                  "numeric"                  "numeric"
\end{verbatim}

\begin{Shaded}
\begin{Highlighting}[]
\FunctionTok{library}\NormalTok{(dplyr)}

\NormalTok{data}\SpecialCharTok{$}\NormalTok{sex }\OtherTok{\textless{}{-}} \FunctionTok{as.factor}\NormalTok{(data}\SpecialCharTok{$}\NormalTok{sex)}
\NormalTok{data}\SpecialCharTok{$}\NormalTok{measurement\_date }\OtherTok{\textless{}{-}} \FunctionTok{as.Date}\NormalTok{(data}\SpecialCharTok{$}\NormalTok{measurement\_date)}

\ControlFlowTok{if}\NormalTok{ (}\StringTok{"participant\_id"} \SpecialCharTok{\%in\%} \FunctionTok{colnames}\NormalTok{(data)) \{}
\NormalTok{  data }\OtherTok{\textless{}{-}}\NormalTok{ data }\SpecialCharTok{\%\textgreater{}\%}\NormalTok{ dplyr}\SpecialCharTok{::}\FunctionTok{select}\NormalTok{(}\SpecialCharTok{{-}}\NormalTok{participant\_id)}
\NormalTok{\}}

\FunctionTok{head}\NormalTok{(data)}
\end{Highlighting}
\end{Shaded}

\begin{verbatim}
## # A tibble: 6 x 9
##   sex     age measurement_date   bmi percent_body_fat hand_grip_strength_kg
##   <fct> <dbl> <date>           <dbl>            <dbl>                 <dbl>
## 1 M        46 2019-01-19        26.5             45                    47  
## 2 M        25 2018-03-16        28.8             45                    28  
## 3 M        53 2017-06-02        25.7             45                    37.7
## 4 M        46 2018-02-11        28               45                    48.9
## 5 F        32 2018-01-25        25               57.7                  20  
## 6 F        28 2018-11-06        24.7             60                    21.2
## # i 3 more variables: sit_and_reach_cm <dbl>, sit_ups_count <dbl>,
## #   vo2_estimate_ml_per_kg_min <dbl>
\end{verbatim}

\begin{Shaded}
\begin{Highlighting}[]
\NormalTok{num\_vars }\OtherTok{\textless{}{-}} \FunctionTok{c}\NormalTok{(}
  \StringTok{"age"}\NormalTok{,}
  \StringTok{"bmi"}\NormalTok{,}
  \StringTok{"percent\_body\_fat"}\NormalTok{,}
  \StringTok{"hand\_grip\_strength\_kg"}\NormalTok{,}
  \StringTok{"sit\_and\_reach\_cm"}\NormalTok{,}
  \StringTok{"sit\_ups\_count"}\NormalTok{,}
  \StringTok{"vo2\_estimate\_ml\_per\_kg\_min"}
\NormalTok{)}

\NormalTok{data\_numeric }\OtherTok{\textless{}{-}}\NormalTok{ data }\SpecialCharTok{\%\textgreater{}\%} \FunctionTok{select}\NormalTok{(}\FunctionTok{all\_of}\NormalTok{(num\_vars))}

\FunctionTok{summary}\NormalTok{(data\_numeric)}
\end{Highlighting}
\end{Shaded}

\begin{verbatim}
##       age             bmi        percent_body_fat hand_grip_strength_kg
##  Min.   :19.00   Min.   :15.00   Min.   :26.50    Min.   : 7.40        
##  1st Qu.:30.00   1st Qu.:21.80   1st Qu.:45.00    1st Qu.:27.88        
##  Median :41.00   Median :24.00   Median :45.00    Median :36.40        
##  Mean   :41.54   Mean   :23.98   Mean   :47.80    Mean   :35.96        
##  3rd Qu.:53.00   3rd Qu.:26.10   3rd Qu.:53.33    3rd Qu.:43.70        
##  Max.   :64.00   Max.   :37.10   Max.   :60.00    Max.   :68.00        
##  sit_and_reach_cm  sit_ups_count   vo2_estimate_ml_per_kg_min
##  Min.   :-14.900   Min.   : 5.00   Min.   :10.00             
##  1st Qu.:  3.700   1st Qu.:15.00   1st Qu.:31.50             
##  Median :  7.900   Median :17.00   Median :36.40             
##  Mean   :  7.887   Mean   :17.62   Mean   :36.28             
##  3rd Qu.: 11.900   3rd Qu.:20.00   3rd Qu.:41.20             
##  Max.   : 34.400   Max.   :33.00   Max.   :60.10
\end{verbatim}

\begin{center}\rule{0.5\linewidth}{0.5pt}\end{center}

\section{Analyse Univariée}\label{analyse-univariuxe9e}

\begin{Shaded}
\begin{Highlighting}[]
\FunctionTok{summary}\NormalTok{(data)}
\end{Highlighting}
\end{Shaded}

\begin{verbatim}
##  sex           age        measurement_date          bmi        percent_body_fat
##  F: 645   Min.   :19.00   Min.   :2015-01-01   Min.   :15.00   Min.   :26.50   
##  M:1355   1st Qu.:30.00   1st Qu.:2016-04-07   1st Qu.:21.80   1st Qu.:45.00   
##           Median :41.00   Median :2017-06-29   Median :24.00   Median :45.00   
##           Mean   :41.54   Mean   :2017-07-01   Mean   :23.98   Mean   :47.80   
##           3rd Qu.:53.00   3rd Qu.:2018-09-18   3rd Qu.:26.10   3rd Qu.:53.33   
##           Max.   :64.00   Max.   :2019-12-31   Max.   :37.10   Max.   :60.00   
##  hand_grip_strength_kg sit_and_reach_cm  sit_ups_count  
##  Min.   : 7.40         Min.   :-14.900   Min.   : 5.00  
##  1st Qu.:27.88         1st Qu.:  3.700   1st Qu.:15.00  
##  Median :36.40         Median :  7.900   Median :17.00  
##  Mean   :35.96         Mean   :  7.887   Mean   :17.62  
##  3rd Qu.:43.70         3rd Qu.: 11.900   3rd Qu.:20.00  
##  Max.   :68.00         Max.   : 34.400   Max.   :33.00  
##  vo2_estimate_ml_per_kg_min
##  Min.   :10.00             
##  1st Qu.:31.50             
##  Median :36.40             
##  Mean   :36.28             
##  3rd Qu.:41.20             
##  Max.   :60.10
\end{verbatim}

\subsection{Préparation des
données}\label{pruxe9paration-des-donnuxe9es}

\subsubsection{Vérification des
doublons}\label{vuxe9rification-des-doublons}

\begin{Shaded}
\begin{Highlighting}[]
\FunctionTok{sum}\NormalTok{(}\FunctionTok{duplicated}\NormalTok{(data))}
\end{Highlighting}
\end{Shaded}

\begin{verbatim}
## [1] 0
\end{verbatim}

\subsubsection{Vérification des valeurs
manquantes}\label{vuxe9rification-des-valeurs-manquantes}

\begin{Shaded}
\begin{Highlighting}[]
\FunctionTok{colSums}\NormalTok{(}\FunctionTok{is.na}\NormalTok{(data))}
\end{Highlighting}
\end{Shaded}

\begin{verbatim}
##                        sex                        age 
##                          0                          0 
##           measurement_date                        bmi 
##                          0                          0 
##           percent_body_fat      hand_grip_strength_kg 
##                          0                          0 
##           sit_and_reach_cm              sit_ups_count 
##                          0                          0 
## vo2_estimate_ml_per_kg_min 
##                          0
\end{verbatim}

\begin{Shaded}
\begin{Highlighting}[]
\FunctionTok{library}\NormalTok{(naniar)}
\FunctionTok{vis\_miss}\NormalTok{(data)}
\end{Highlighting}
\end{Shaded}

\begin{center}\includegraphics{rapport_ameliore_files/figure-latex/valeurs-manquantes-1} \end{center}

\subsubsection{Détection et correction des valeurs
aberrantes}\label{duxe9tection-et-correction-des-valeurs-aberrantes}

La détection des valeurs aberrantes constitue une étape importante de
l'analyse exploratoire des données, car la présence d'observations
extrêmes peut influencer les statistiques descriptives, les tests
d'hypothèses et les modèles statistiques ultérieurs.

Dans cette étude, les valeurs aberrantes sont identifiées à l'aide d'une
approche graphique basée sur les boxplots. Les observations situées en
dehors des bornes définies par l'intervalle
\([Q_1 - 1.5 \times IQR,\ Q_3 + 1.5 \times IQR]\) sont considérées comme
aberrantes.

\begin{Shaded}
\begin{Highlighting}[]
\CommentTok{\# {-}{-}{-} Fonction de détection {-}{-}{-}}
\NormalTok{detecter\_outliers }\OtherTok{\textless{}{-}} \ControlFlowTok{function}\NormalTok{(variable, nom\_variable) \{}
  \FunctionTok{cat}\NormalTok{(}\StringTok{"}\SpecialCharTok{\textbackslash{}n}\StringTok{====================================}\SpecialCharTok{\textbackslash{}n}\StringTok{"}\NormalTok{)}
  \FunctionTok{cat}\NormalTok{(}\StringTok{"Analyse de :"}\NormalTok{, nom\_variable, }\StringTok{"}\SpecialCharTok{\textbackslash{}n}\StringTok{"}\NormalTok{)}
  \FunctionTok{cat}\NormalTok{(}\StringTok{"====================================}\SpecialCharTok{\textbackslash{}n}\StringTok{"}\NormalTok{)}

\NormalTok{  Q1 }\OtherTok{\textless{}{-}} \FunctionTok{quantile}\NormalTok{(variable, }\FloatTok{0.25}\NormalTok{, }\AttributeTok{na.rm =} \ConstantTok{TRUE}\NormalTok{)}
\NormalTok{  Q3 }\OtherTok{\textless{}{-}} \FunctionTok{quantile}\NormalTok{(variable, }\FloatTok{0.75}\NormalTok{, }\AttributeTok{na.rm =} \ConstantTok{TRUE}\NormalTok{)}
\NormalTok{  IQR\_value }\OtherTok{\textless{}{-}}\NormalTok{ Q3 }\SpecialCharTok{{-}}\NormalTok{ Q1}
\NormalTok{  limite\_basse }\OtherTok{\textless{}{-}}\NormalTok{ Q1 }\SpecialCharTok{{-}} \FloatTok{1.5} \SpecialCharTok{*}\NormalTok{ IQR\_value}
\NormalTok{  limite\_haute }\OtherTok{\textless{}{-}}\NormalTok{ Q3 }\SpecialCharTok{+} \FloatTok{1.5} \SpecialCharTok{*}\NormalTok{ IQR\_value}

\NormalTok{  outliers\_bas  }\OtherTok{\textless{}{-}} \FunctionTok{sum}\NormalTok{(variable }\SpecialCharTok{\textless{}}\NormalTok{ limite\_basse, }\AttributeTok{na.rm =} \ConstantTok{TRUE}\NormalTok{)}
\NormalTok{  outliers\_haut }\OtherTok{\textless{}{-}} \FunctionTok{sum}\NormalTok{(variable }\SpecialCharTok{\textgreater{}}\NormalTok{ limite\_haute, }\AttributeTok{na.rm =} \ConstantTok{TRUE}\NormalTok{)}
\NormalTok{  total\_outliers }\OtherTok{\textless{}{-}}\NormalTok{ outliers\_bas }\SpecialCharTok{+}\NormalTok{ outliers\_haut}

  \FunctionTok{cat}\NormalTok{(}\StringTok{"Outliers bas  :"}\NormalTok{, outliers\_bas,  }\StringTok{"}\SpecialCharTok{\textbackslash{}n}\StringTok{"}\NormalTok{)}
  \FunctionTok{cat}\NormalTok{(}\StringTok{"Outliers haut :"}\NormalTok{, outliers\_haut, }\StringTok{"}\SpecialCharTok{\textbackslash{}n}\StringTok{"}\NormalTok{)}
  \FunctionTok{cat}\NormalTok{(}\StringTok{"Total         :"}\NormalTok{, total\_outliers, }\StringTok{"}\SpecialCharTok{\textbackslash{}n}\StringTok{"}\NormalTok{)}
  \FunctionTok{cat}\NormalTok{(}\StringTok{"Limite basse  :"}\NormalTok{, }\FunctionTok{round}\NormalTok{(limite\_basse, }\DecValTok{2}\NormalTok{), }\StringTok{"}\SpecialCharTok{\textbackslash{}n}\StringTok{"}\NormalTok{)}
  \FunctionTok{cat}\NormalTok{(}\StringTok{"Limite haute  :"}\NormalTok{, }\FunctionTok{round}\NormalTok{(limite\_haute, }\DecValTok{2}\NormalTok{), }\StringTok{"}\SpecialCharTok{\textbackslash{}n}\StringTok{"}\NormalTok{)}

  \FunctionTok{return}\NormalTok{(}\FunctionTok{list}\NormalTok{(}\AttributeTok{limite\_basse =}\NormalTok{ limite\_basse, }\AttributeTok{limite\_haute =}\NormalTok{ limite\_haute))}
\NormalTok{\}}

\CommentTok{\# {-}{-}{-} Fonction de correction (Winsorisation) {-}{-}{-}}
\NormalTok{corriger\_outliers }\OtherTok{\textless{}{-}} \ControlFlowTok{function}\NormalTok{(variable) \{}
\NormalTok{  Q1 }\OtherTok{\textless{}{-}} \FunctionTok{quantile}\NormalTok{(variable, }\FloatTok{0.25}\NormalTok{, }\AttributeTok{na.rm =} \ConstantTok{TRUE}\NormalTok{)}
\NormalTok{  Q3 }\OtherTok{\textless{}{-}} \FunctionTok{quantile}\NormalTok{(variable, }\FloatTok{0.75}\NormalTok{, }\AttributeTok{na.rm =} \ConstantTok{TRUE}\NormalTok{)}
\NormalTok{  IQR\_value }\OtherTok{\textless{}{-}}\NormalTok{ Q3 }\SpecialCharTok{{-}}\NormalTok{ Q1}
\NormalTok{  limite\_basse }\OtherTok{\textless{}{-}}\NormalTok{ Q1 }\SpecialCharTok{{-}} \FloatTok{1.5} \SpecialCharTok{*}\NormalTok{ IQR\_value}
\NormalTok{  limite\_haute }\OtherTok{\textless{}{-}}\NormalTok{ Q3 }\SpecialCharTok{+} \FloatTok{1.5} \SpecialCharTok{*}\NormalTok{ IQR\_value}

\NormalTok{  variable\_corrigee }\OtherTok{\textless{}{-}}\NormalTok{ variable}
\NormalTok{  variable\_corrigee[variable }\SpecialCharTok{\textless{}}\NormalTok{ limite\_basse] }\OtherTok{\textless{}{-}}\NormalTok{ limite\_basse}
\NormalTok{  variable\_corrigee[variable }\SpecialCharTok{\textgreater{}}\NormalTok{ limite\_haute] }\OtherTok{\textless{}{-}}\NormalTok{ limite\_haute}

  \FunctionTok{return}\NormalTok{(variable\_corrigee)}
\NormalTok{\}}

\CommentTok{\# {-}{-}{-} Analyse des valeurs aberrantes {-}{-}{-}}
\FunctionTok{cat}\NormalTok{(}\StringTok{"}\SpecialCharTok{\textbackslash{}n}\StringTok{ANALYSE DES VALEURS ABERRANTES}\SpecialCharTok{\textbackslash{}n}\StringTok{"}\NormalTok{)}
\end{Highlighting}
\end{Shaded}

\begin{verbatim}
## 
## ANALYSE DES VALEURS ABERRANTES
\end{verbatim}

\begin{Shaded}
\begin{Highlighting}[]
\FunctionTok{cat}\NormalTok{(}\StringTok{"====================================}\SpecialCharTok{\textbackslash{}n}\StringTok{"}\NormalTok{)}
\end{Highlighting}
\end{Shaded}

\begin{verbatim}
## ====================================
\end{verbatim}

\begin{Shaded}
\begin{Highlighting}[]
\NormalTok{variables\_numeriques }\OtherTok{\textless{}{-}} \FunctionTok{c}\NormalTok{(}
  \StringTok{"age"}\NormalTok{, }\StringTok{"bmi"}\NormalTok{, }\StringTok{"percent\_body\_fat"}\NormalTok{,}
  \StringTok{"hand\_grip\_strength\_kg"}\NormalTok{, }\StringTok{"sit\_and\_reach\_cm"}\NormalTok{,}
  \StringTok{"sit\_ups\_count"}\NormalTok{, }\StringTok{"vo2\_estimate\_ml\_per\_kg\_min"}
\NormalTok{)}

\ControlFlowTok{for}\NormalTok{ (var }\ControlFlowTok{in}\NormalTok{ variables\_numeriques) \{}
  \ControlFlowTok{if}\NormalTok{ (var }\SpecialCharTok{\%in\%} \FunctionTok{names}\NormalTok{(data)) \{}
    \FunctionTok{detecter\_outliers}\NormalTok{(data[[var]], var)}
\NormalTok{  \}}
\NormalTok{\}}
\end{Highlighting}
\end{Shaded}

\begin{verbatim}
## 
## ====================================
## Analyse de : age 
## ====================================
## Outliers bas  : 0 
## Outliers haut : 0 
## Total         : 0 
## Limite basse  : -4.5 
## Limite haute  : 87.5 
## 
## ====================================
## Analyse de : bmi 
## ====================================
## Outliers bas  : 8 
## Outliers haut : 5 
## Total         : 13 
## Limite basse  : 15.35 
## Limite haute  : 32.55 
## 
## ====================================
## Analyse de : percent_body_fat 
## ====================================
## Outliers bas  : 14 
## Outliers haut : 0 
## Total         : 14 
## Limite basse  : 32.51 
## Limite haute  : 65.81 
## 
## ====================================
## Analyse de : hand_grip_strength_kg 
## ====================================
## Outliers bas  : 0 
## Outliers haut : 1 
## Total         : 1 
## Limite basse  : 4.14 
## Limite haute  : 67.44 
## 
## ====================================
## Analyse de : sit_and_reach_cm 
## ====================================
## Outliers bas  : 10 
## Outliers haut : 7 
## Total         : 17 
## Limite basse  : -8.6 
## Limite haute  : 24.2 
## 
## ====================================
## Analyse de : sit_ups_count 
## ====================================
## Outliers bas  : 15 
## Outliers haut : 32 
## Total         : 47 
## Limite basse  : 7.5 
## Limite haute  : 27.5 
## 
## ====================================
## Analyse de : vo2_estimate_ml_per_kg_min 
## ====================================
## Outliers bas  : 5 
## Outliers haut : 7 
## Total         : 12 
## Limite basse  : 16.95 
## Limite haute  : 55.75
\end{verbatim}

\begin{Shaded}
\begin{Highlighting}[]
\CommentTok{\# {-}{-}{-} Correction des outliers {-}{-}{-}}
\FunctionTok{cat}\NormalTok{(}\StringTok{"}\SpecialCharTok{\textbackslash{}n}\StringTok{CORRECTION DES VALEURS ABERRANTES}\SpecialCharTok{\textbackslash{}n}\StringTok{"}\NormalTok{)}
\end{Highlighting}
\end{Shaded}

\begin{verbatim}
## 
## CORRECTION DES VALEURS ABERRANTES
\end{verbatim}

\begin{Shaded}
\begin{Highlighting}[]
\FunctionTok{cat}\NormalTok{(}\StringTok{"====================================}\SpecialCharTok{\textbackslash{}n}\StringTok{"}\NormalTok{)}
\end{Highlighting}
\end{Shaded}

\begin{verbatim}
## ====================================
\end{verbatim}

\begin{Shaded}
\begin{Highlighting}[]
\NormalTok{health\_clean }\OtherTok{\textless{}{-}}\NormalTok{ data}
\NormalTok{variables\_a\_corriger }\OtherTok{\textless{}{-}} \FunctionTok{c}\NormalTok{(}\StringTok{"bmi"}\NormalTok{, }\StringTok{"percent\_body\_fat"}\NormalTok{, }\StringTok{"hand\_grip\_strength\_kg"}\NormalTok{)}

\ControlFlowTok{for}\NormalTok{ (var }\ControlFlowTok{in}\NormalTok{ variables\_a\_corriger) \{}
  \ControlFlowTok{if}\NormalTok{ (var }\SpecialCharTok{\%in\%} \FunctionTok{names}\NormalTok{(data)) \{}
\NormalTok{    avant }\OtherTok{\textless{}{-}}\NormalTok{ data[[var]]}
\NormalTok{    apres }\OtherTok{\textless{}{-}} \FunctionTok{corriger\_outliers}\NormalTok{(avant)}
\NormalTok{    health\_clean[[var]] }\OtherTok{\textless{}{-}}\NormalTok{ apres}
\NormalTok{    modifications }\OtherTok{\textless{}{-}} \FunctionTok{sum}\NormalTok{(avant }\SpecialCharTok{!=}\NormalTok{ apres, }\AttributeTok{na.rm =} \ConstantTok{TRUE}\NormalTok{)}
    \FunctionTok{cat}\NormalTok{(var, }\StringTok{":"}\NormalTok{, modifications, }\StringTok{"valeurs corrigées}\SpecialCharTok{\textbackslash{}n}\StringTok{"}\NormalTok{)}
\NormalTok{  \}}
\NormalTok{\}}
\end{Highlighting}
\end{Shaded}

\begin{verbatim}
## bmi : 13 valeurs corrigées
## percent_body_fat : 14 valeurs corrigées
## hand_grip_strength_kg : 1 valeurs corrigées
\end{verbatim}

\begin{Shaded}
\begin{Highlighting}[]
\CommentTok{\# {-}{-}{-} Comparaison avant / après {-}{-}{-}}
\FunctionTok{cat}\NormalTok{(}\StringTok{"}\SpecialCharTok{\textbackslash{}n}\StringTok{COMPARAISON AVANT / APRÈS}\SpecialCharTok{\textbackslash{}n}\StringTok{"}\NormalTok{)}
\end{Highlighting}
\end{Shaded}

\begin{verbatim}
## 
## COMPARAISON AVANT / APRÈS
\end{verbatim}

\begin{Shaded}
\begin{Highlighting}[]
\FunctionTok{cat}\NormalTok{(}\StringTok{"====================================}\SpecialCharTok{\textbackslash{}n}\StringTok{"}\NormalTok{)}
\end{Highlighting}
\end{Shaded}

\begin{verbatim}
## ====================================
\end{verbatim}

\begin{Shaded}
\begin{Highlighting}[]
\NormalTok{comparer\_stats }\OtherTok{\textless{}{-}} \ControlFlowTok{function}\NormalTok{(avant, apres, nom) \{}
  \FunctionTok{cat}\NormalTok{(}\StringTok{"}\SpecialCharTok{\textbackslash{}n}\StringTok{Variable :"}\NormalTok{, nom, }\StringTok{"}\SpecialCharTok{\textbackslash{}n}\StringTok{"}\NormalTok{)}
  \FunctionTok{cat}\NormalTok{(}\StringTok{"AVANT {-}\textgreater{} Min :"}\NormalTok{, }\FunctionTok{round}\NormalTok{(}\FunctionTok{min}\NormalTok{(avant, }\AttributeTok{na.rm =} \ConstantTok{TRUE}\NormalTok{), }\DecValTok{2}\NormalTok{),}
      \StringTok{"| Max :"}\NormalTok{, }\FunctionTok{round}\NormalTok{(}\FunctionTok{max}\NormalTok{(avant, }\AttributeTok{na.rm =} \ConstantTok{TRUE}\NormalTok{), }\DecValTok{2}\NormalTok{),}
      \StringTok{"| Moyenne :"}\NormalTok{, }\FunctionTok{round}\NormalTok{(}\FunctionTok{mean}\NormalTok{(avant, }\AttributeTok{na.rm =} \ConstantTok{TRUE}\NormalTok{), }\DecValTok{2}\NormalTok{), }\StringTok{"}\SpecialCharTok{\textbackslash{}n}\StringTok{"}\NormalTok{)}
  \FunctionTok{cat}\NormalTok{(}\StringTok{"APRÈS {-}\textgreater{} Min :"}\NormalTok{, }\FunctionTok{round}\NormalTok{(}\FunctionTok{min}\NormalTok{(apres, }\AttributeTok{na.rm =} \ConstantTok{TRUE}\NormalTok{), }\DecValTok{2}\NormalTok{),}
      \StringTok{"| Max :"}\NormalTok{, }\FunctionTok{round}\NormalTok{(}\FunctionTok{max}\NormalTok{(apres, }\AttributeTok{na.rm =} \ConstantTok{TRUE}\NormalTok{), }\DecValTok{2}\NormalTok{),}
      \StringTok{"| Moyenne :"}\NormalTok{, }\FunctionTok{round}\NormalTok{(}\FunctionTok{mean}\NormalTok{(apres, }\AttributeTok{na.rm =} \ConstantTok{TRUE}\NormalTok{), }\DecValTok{2}\NormalTok{), }\StringTok{"}\SpecialCharTok{\textbackslash{}n}\StringTok{"}\NormalTok{)}
\NormalTok{\}}

\ControlFlowTok{for}\NormalTok{ (var }\ControlFlowTok{in}\NormalTok{ variables\_a\_corriger) \{}
  \ControlFlowTok{if}\NormalTok{ (var }\SpecialCharTok{\%in\%} \FunctionTok{names}\NormalTok{(data)) \{}
    \FunctionTok{comparer\_stats}\NormalTok{(data[[var]], health\_clean[[var]], var)}
\NormalTok{  \}}
\NormalTok{\}}
\end{Highlighting}
\end{Shaded}

\begin{verbatim}
## 
## Variable : bmi 
## AVANT -> Min : 15 | Max : 37.1 | Moyenne : 23.98 
## APRÈS -> Min : 15.35 | Max : 32.55 | Moyenne : 23.98 
## 
## Variable : percent_body_fat 
## AVANT -> Min : 26.5 | Max : 60 | Moyenne : 47.8 
## APRÈS -> Min : 32.51 | Max : 60 | Moyenne : 47.81 
## 
## Variable : hand_grip_strength_kg 
## AVANT -> Min : 7.4 | Max : 68 | Moyenne : 35.96 
## APRÈS -> Min : 7.4 | Max : 67.44 | Moyenne : 35.96
\end{verbatim}

\begin{Shaded}
\begin{Highlighting}[]
\CommentTok{\# {-}{-}{-} Visualisation {-}{-}{-}}
\NormalTok{data\_compare }\OtherTok{\textless{}{-}} \FunctionTok{data.frame}\NormalTok{()}

\ControlFlowTok{for}\NormalTok{ (var }\ControlFlowTok{in}\NormalTok{ variables\_a\_corriger) \{}
\NormalTok{  temp }\OtherTok{\textless{}{-}} \FunctionTok{data.frame}\NormalTok{(}
    \AttributeTok{Valeur   =} \FunctionTok{c}\NormalTok{(data[[var]], health\_clean[[var]]),}
    \AttributeTok{Variable =}\NormalTok{ var,}
    \AttributeTok{Statut   =} \FunctionTok{rep}\NormalTok{(}\FunctionTok{c}\NormalTok{(}\StringTok{"Avant"}\NormalTok{, }\StringTok{"Après"}\NormalTok{), }\AttributeTok{each =} \FunctionTok{nrow}\NormalTok{(data))}
\NormalTok{  )}
\NormalTok{  data\_compare }\OtherTok{\textless{}{-}} \FunctionTok{rbind}\NormalTok{(data\_compare, temp)}
\NormalTok{\}}

\FunctionTok{ggplot}\NormalTok{(data\_compare, }\FunctionTok{aes}\NormalTok{(}\AttributeTok{x =}\NormalTok{ Variable, }\AttributeTok{y =}\NormalTok{ Valeur, }\AttributeTok{fill =}\NormalTok{ Statut)) }\SpecialCharTok{+}
  \FunctionTok{geom\_boxplot}\NormalTok{() }\SpecialCharTok{+}
  \FunctionTok{labs}\NormalTok{(}
    \AttributeTok{title =} \StringTok{"Comparaison avant / après correction des valeurs aberrantes"}\NormalTok{,}
    \AttributeTok{x =} \StringTok{"Variables"}\NormalTok{,}
    \AttributeTok{y =} \StringTok{"Valeurs"}
\NormalTok{  ) }\SpecialCharTok{+}
  \FunctionTok{theme\_minimal}\NormalTok{()}
\end{Highlighting}
\end{Shaded}

\begin{center}\includegraphics{rapport_ameliore_files/figure-latex/outliers-1} \end{center}

\begin{Shaded}
\begin{Highlighting}[]
\CommentTok{\# {-}{-}{-} Sauvegarde {-}{-}{-}}
\FunctionTok{write.csv}\NormalTok{(health\_clean, }\StringTok{"health\_dataset\_clean.csv"}\NormalTok{, }\AttributeTok{row.names =} \ConstantTok{FALSE}\NormalTok{)}

\FunctionTok{cat}\NormalTok{(}\StringTok{"}\SpecialCharTok{\textbackslash{}n}\StringTok{RAPPORT FINAL}\SpecialCharTok{\textbackslash{}n}\StringTok{"}\NormalTok{)}
\end{Highlighting}
\end{Shaded}

\begin{verbatim}
## 
## RAPPORT FINAL
\end{verbatim}

\begin{Shaded}
\begin{Highlighting}[]
\FunctionTok{cat}\NormalTok{(}\StringTok{"====================================}\SpecialCharTok{\textbackslash{}n}\StringTok{"}\NormalTok{)}
\end{Highlighting}
\end{Shaded}

\begin{verbatim}
## ====================================
\end{verbatim}

\begin{Shaded}
\begin{Highlighting}[]
\FunctionTok{cat}\NormalTok{(}\StringTok{"Observations totales :"}\NormalTok{, }\FunctionTok{nrow}\NormalTok{(health\_clean), }\StringTok{"}\SpecialCharTok{\textbackslash{}n}\StringTok{"}\NormalTok{)}
\end{Highlighting}
\end{Shaded}

\begin{verbatim}
## Observations totales : 2000
\end{verbatim}

\begin{Shaded}
\begin{Highlighting}[]
\FunctionTok{cat}\NormalTok{(}\StringTok{"Variables nettoyées  :"}\NormalTok{, }\FunctionTok{length}\NormalTok{(variables\_numeriques), }\StringTok{"}\SpecialCharTok{\textbackslash{}n}\StringTok{"}\NormalTok{)}
\end{Highlighting}
\end{Shaded}

\begin{verbatim}
## Variables nettoyées  : 7
\end{verbatim}

\begin{Shaded}
\begin{Highlighting}[]
\FunctionTok{cat}\NormalTok{(}\StringTok{"Fichier sauvegardé   : health\_dataset\_clean.csv}\SpecialCharTok{\textbackslash{}n}\StringTok{"}\NormalTok{)}
\end{Highlighting}
\end{Shaded}

\begin{verbatim}
## Fichier sauvegardé   : health_dataset_clean.csv
\end{verbatim}

\begin{Shaded}
\begin{Highlighting}[]
\FunctionTok{cat}\NormalTok{(}\StringTok{"Nettoyage terminé avec succès !}\SpecialCharTok{\textbackslash{}n\textbackslash{}n}\StringTok{"}\NormalTok{)}
\end{Highlighting}
\end{Shaded}

\begin{verbatim}
## Nettoyage terminé avec succès !
\end{verbatim}

\begin{Shaded}
\begin{Highlighting}[]
\FunctionTok{print}\NormalTok{(}\FunctionTok{head}\NormalTok{(health\_clean))}
\end{Highlighting}
\end{Shaded}

\begin{verbatim}
## # A tibble: 6 x 9
##   sex     age measurement_date   bmi percent_body_fat hand_grip_strength_kg
##   <fct> <dbl> <date>           <dbl>            <dbl>                 <dbl>
## 1 M        46 2019-01-19        26.5             45                    47  
## 2 M        25 2018-03-16        28.8             45                    28  
## 3 M        53 2017-06-02        25.7             45                    37.7
## 4 M        46 2018-02-11        28               45                    48.9
## 5 F        32 2018-01-25        25               57.7                  20  
## 6 F        28 2018-11-06        24.7             60                    21.2
## # i 3 more variables: sit_and_reach_cm <dbl>, sit_ups_count <dbl>,
## #   vo2_estimate_ml_per_kg_min <dbl>
\end{verbatim}

\subsection{Analyse de la distribution des
variables}\label{analyse-de-la-distribution-des-variables}

L'étude de la distribution des variables quantitatives constitue une
étape clé de l'analyse exploratoire. Elle permet d'évaluer visuellement
la forme des distributions, d'identifier d'éventuelles asymétries et
d'anticiper le choix des tests statistiques appropriés.

Pour chaque variable numérique, un histogramme accompagné de sa courbe
de densité est tracé, ainsi qu'un QQ-plot et un test de Shapiro-Wilk.

\begin{Shaded}
\begin{Highlighting}[]
\FunctionTok{library}\NormalTok{(ggplot2)}

\NormalTok{target\_vars }\OtherTok{\textless{}{-}} \FunctionTok{c}\NormalTok{(}
  \StringTok{"hand\_grip\_strength\_kg"}\NormalTok{,}
  \StringTok{"sit\_and\_reach\_cm"}\NormalTok{,}
  \StringTok{"sit\_ups\_count"}\NormalTok{,}
  \StringTok{"vo2\_estimate\_ml\_per\_kg\_min"}
\NormalTok{)}

\ControlFlowTok{for}\NormalTok{ (var }\ControlFlowTok{in}\NormalTok{ target\_vars) \{}
  \FunctionTok{cat}\NormalTok{(}\StringTok{"{-}{-}{-}{-}{-}{-}{-}{-}{-}{-}{-}{-}{-}{-}{-}{-}{-}{-}{-}{-}{-}{-}{-}{-}{-}{-}{-}{-}{-}{-}{-}{-}{-}{-}{-}{-}{-}{-}{-}{-}{-}{-}{-}{-}{-}{-}{-}{-}{-}{-}}\SpecialCharTok{\textbackslash{}n}\StringTok{"}\NormalTok{)}
  \FunctionTok{cat}\NormalTok{(}\StringTok{"Variable :"}\NormalTok{, var, }\StringTok{"}\SpecialCharTok{\textbackslash{}n}\StringTok{"}\NormalTok{)}
  \FunctionTok{cat}\NormalTok{(}\StringTok{"{-}{-}{-}{-}{-}{-}{-}{-}{-}{-}{-}{-}{-}{-}{-}{-}{-}{-}{-}{-}{-}{-}{-}{-}{-}{-}{-}{-}{-}{-}{-}{-}{-}{-}{-}{-}{-}{-}{-}{-}{-}{-}{-}{-}{-}{-}{-}{-}{-}{-}}\SpecialCharTok{\textbackslash{}n\textbackslash{}n}\StringTok{"}\NormalTok{)}

\NormalTok{  x }\OtherTok{\textless{}{-}}\NormalTok{ health\_clean[[var]]}
  \FunctionTok{print}\NormalTok{(}\FunctionTok{summary}\NormalTok{(x))}

\NormalTok{  mean\_x }\OtherTok{\textless{}{-}} \FunctionTok{mean}\NormalTok{(x, }\AttributeTok{na.rm =} \ConstantTok{TRUE}\NormalTok{)}
\NormalTok{  sd\_x   }\OtherTok{\textless{}{-}} \FunctionTok{sd}\NormalTok{(x, }\AttributeTok{na.rm =} \ConstantTok{TRUE}\NormalTok{)}
\NormalTok{  df\_var }\OtherTok{\textless{}{-}} \FunctionTok{data.frame}\NormalTok{(}\AttributeTok{x =}\NormalTok{ x)}

  \CommentTok{\# Histogramme + densité}
\NormalTok{  p1 }\OtherTok{\textless{}{-}} \FunctionTok{ggplot}\NormalTok{(df\_var, }\FunctionTok{aes}\NormalTok{(x)) }\SpecialCharTok{+}
    \FunctionTok{geom\_histogram}\NormalTok{(}\FunctionTok{aes}\NormalTok{(}\AttributeTok{y =}\NormalTok{ ..density..), }\AttributeTok{bins =} \DecValTok{30}\NormalTok{,}
                   \AttributeTok{fill =} \StringTok{"lightblue"}\NormalTok{, }\AttributeTok{alpha =} \FloatTok{0.6}\NormalTok{) }\SpecialCharTok{+}
    \FunctionTok{geom\_density}\NormalTok{(}\AttributeTok{color =} \StringTok{"red"}\NormalTok{, }\AttributeTok{linewidth =} \FloatTok{1.2}\NormalTok{) }\SpecialCharTok{+}
    \FunctionTok{stat\_function}\NormalTok{(}\AttributeTok{fun =}\NormalTok{ dnorm,}
                  \AttributeTok{args =} \FunctionTok{list}\NormalTok{(}\AttributeTok{mean =}\NormalTok{ mean\_x, }\AttributeTok{sd =}\NormalTok{ sd\_x),}
                  \AttributeTok{color =} \StringTok{"darkgreen"}\NormalTok{, }\AttributeTok{linewidth =} \FloatTok{1.2}\NormalTok{) }\SpecialCharTok{+}
    \FunctionTok{theme\_minimal}\NormalTok{() }\SpecialCharTok{+}
    \FunctionTok{labs}\NormalTok{(}\AttributeTok{title =} \FunctionTok{paste}\NormalTok{(}\StringTok{"Distribution —"}\NormalTok{, var), }\AttributeTok{x =}\NormalTok{ var, }\AttributeTok{y =} \StringTok{"Densité"}\NormalTok{)}

  \FunctionTok{print}\NormalTok{(p1)}

  \CommentTok{\# QQ Plot}
\NormalTok{  p2 }\OtherTok{\textless{}{-}} \FunctionTok{ggplot}\NormalTok{(df\_var, }\FunctionTok{aes}\NormalTok{(}\AttributeTok{sample =}\NormalTok{ x)) }\SpecialCharTok{+}
    \FunctionTok{stat\_qq}\NormalTok{(}\AttributeTok{color =} \StringTok{"blue"}\NormalTok{) }\SpecialCharTok{+}
    \FunctionTok{stat\_qq\_line}\NormalTok{(}\AttributeTok{color =} \StringTok{"red"}\NormalTok{) }\SpecialCharTok{+}
    \FunctionTok{theme\_minimal}\NormalTok{() }\SpecialCharTok{+}
    \FunctionTok{labs}\NormalTok{(}\AttributeTok{title =} \FunctionTok{paste}\NormalTok{(}\StringTok{"QQ Plot —"}\NormalTok{, var))}

  \FunctionTok{print}\NormalTok{(p2)}

  \CommentTok{\# Test de Shapiro{-}Wilk}
  \ControlFlowTok{if}\NormalTok{ (}\FunctionTok{length}\NormalTok{(}\FunctionTok{na.omit}\NormalTok{(x)) }\SpecialCharTok{\textless{}=} \DecValTok{5000}\NormalTok{) \{}
\NormalTok{    sh }\OtherTok{\textless{}{-}} \FunctionTok{shapiro.test}\NormalTok{(x)}
    \FunctionTok{print}\NormalTok{(sh)}
\NormalTok{  \} }\ControlFlowTok{else}\NormalTok{ \{}
    \FunctionTok{cat}\NormalTok{(}\StringTok{"Test de Shapiro{-}Wilk non applicable (n \textgreater{} 5000 observations)}\SpecialCharTok{\textbackslash{}n}\StringTok{"}\NormalTok{)}
\NormalTok{  \}}

  \FunctionTok{cat}\NormalTok{(}\StringTok{"}\SpecialCharTok{\textbackslash{}n\textbackslash{}n}\StringTok{"}\NormalTok{)}
\NormalTok{\}}
\end{Highlighting}
\end{Shaded}

\begin{verbatim}
## --------------------------------------------------
## Variable : hand_grip_strength_kg 
## --------------------------------------------------
## 
##    Min. 1st Qu.  Median    Mean 3rd Qu.    Max. 
##    7.40   27.88   36.40   35.96   43.70   67.44
\end{verbatim}

\begin{center}\includegraphics{rapport_ameliore_files/figure-latex/distribution-1} \end{center}

\begin{center}\includegraphics{rapport_ameliore_files/figure-latex/distribution-2} \end{center}

\begin{verbatim}
## 
##  Shapiro-Wilk normality test
## 
## data:  x
## W = 0.99336, p-value = 7.622e-08
## 
## 
## 
## --------------------------------------------------
## Variable : sit_and_reach_cm 
## --------------------------------------------------
## 
##    Min. 1st Qu.  Median    Mean 3rd Qu.    Max. 
## -14.900   3.700   7.900   7.887  11.900  34.400
\end{verbatim}

\begin{center}\includegraphics{rapport_ameliore_files/figure-latex/distribution-3} \end{center}

\begin{center}\includegraphics{rapport_ameliore_files/figure-latex/distribution-4} \end{center}

\begin{verbatim}
## 
##  Shapiro-Wilk normality test
## 
## data:  x
## W = 0.99911, p-value = 0.441
## 
## 
## 
## --------------------------------------------------
## Variable : sit_ups_count 
## --------------------------------------------------
## 
##    Min. 1st Qu.  Median    Mean 3rd Qu.    Max. 
##    5.00   15.00   17.00   17.62   20.00   33.00
\end{verbatim}

\begin{center}\includegraphics{rapport_ameliore_files/figure-latex/distribution-5} \end{center}

\begin{center}\includegraphics{rapport_ameliore_files/figure-latex/distribution-6} \end{center}

\begin{verbatim}
## 
##  Shapiro-Wilk normality test
## 
## data:  x
## W = 0.99218, p-value = 7.357e-09
## 
## 
## 
## --------------------------------------------------
## Variable : vo2_estimate_ml_per_kg_min 
## --------------------------------------------------
## 
##    Min. 1st Qu.  Median    Mean 3rd Qu.    Max. 
##   10.00   31.50   36.40   36.28   41.20   60.10
\end{verbatim}

\begin{center}\includegraphics{rapport_ameliore_files/figure-latex/distribution-7} \end{center}

\begin{center}\includegraphics{rapport_ameliore_files/figure-latex/distribution-8} \end{center}

\begin{verbatim}
## 
##  Shapiro-Wilk normality test
## 
## data:  x
## W = 0.99956, p-value = 0.9484
\end{verbatim}

\subsection{Standardisation des
données}\label{standardisation-des-donnuxe9es}

Avant de réaliser l'Analyse en Composantes Principales (ACP), il est
nécessaire de standardiser les variables quantitatives. En effet, les
variables étudiées ne sont pas exprimées dans les mêmes unités (kg, cm,
années, etc.), ce qui pourrait biaiser l'analyse en donnant plus de
poids aux variables ayant de grandes valeurs numériques.

La standardisation transforme les variables afin qu'elles aient une
\textbf{moyenne nulle} et un \textbf{écart-type égal à 1}, garantissant
une contribution équivalente de chaque variable à l'ACP.

\begin{Shaded}
\begin{Highlighting}[]
\FunctionTok{library}\NormalTok{(FactoMineR)}
\FunctionTok{library}\NormalTok{(factoextra)}

\NormalTok{variables\_quantitatives }\OtherTok{\textless{}{-}}\NormalTok{ num\_vars}

\NormalTok{data\_numeric }\OtherTok{\textless{}{-}}\NormalTok{ health\_clean }\SpecialCharTok{\%\textgreater{}\%} \FunctionTok{select}\NormalTok{(}\FunctionTok{all\_of}\NormalTok{(variables\_quantitatives))}
\NormalTok{data\_scaled  }\OtherTok{\textless{}{-}} \FunctionTok{scale}\NormalTok{(data\_numeric)}
\end{Highlighting}
\end{Shaded}

\begin{center}\rule{0.5\linewidth}{0.5pt}\end{center}

\section{Analyse Factorielle --- ACP}\label{analyse-factorielle-acp}

\subsection{Objectif de l'ACP}\label{objectif-de-lacp}

L'Analyse en Composantes Principales (ACP) a été réalisée afin de
\textbf{réduire la dimension} des données physiologiques et
anthropométriques tout en conservant le maximum d'information contenue
dans les variables initiales. Cette méthode permet également
d'identifier les relations existantes entre les variables étudiées ainsi
que les profils des individus selon leurs caractéristiques physiques et
physiologiques.

Les variables quantitatives utilisées sont : âge, IMC, pourcentage de
masse grasse, force de préhension, flexibilité, endurance abdominale
(sit-ups) et VO₂ max.

\begin{Shaded}
\begin{Highlighting}[]
\NormalTok{res.pca }\OtherTok{\textless{}{-}} \FunctionTok{PCA}\NormalTok{(data\_scaled, }\AttributeTok{graph =} \ConstantTok{FALSE}\NormalTok{)}

\NormalTok{eig.val }\OtherTok{\textless{}{-}} \FunctionTok{get\_eigenvalue}\NormalTok{(res.pca)}
\NormalTok{eig.val}
\end{Highlighting}
\end{Shaded}

\begin{verbatim}
##       eigenvalue variance.percent cumulative.variance.percent
## Dim.1  1.9702193        28.145991                    28.14599
## Dim.2  1.2944069        18.491527                    46.63752
## Dim.3  1.1507254        16.438934                    63.07645
## Dim.4  0.9400499        13.429285                    76.50574
## Dim.5  0.7857555        11.225078                    87.73082
## Dim.6  0.5765042         8.235774                    95.96659
## Dim.7  0.2823388         4.033411                   100.00000
\end{verbatim}

\textbf{Interprétation des valeurs propres :}

\begin{itemize}
\tightlist
\item
  La \textbf{1ère composante principale} explique \textbf{28,1 \%} de la
  variance totale.
\item
  La \textbf{2ème composante} explique \textbf{18,5 \%} de la variance
  supplémentaire.
\item
  Les \textbf{deux premiers axes} résument ensemble \textbf{46,6 \%} de
  l'information totale.
\item
  Les \textbf{trois premières composantes} expliquent \textbf{63,1 \%}
  de la variance totale.
\end{itemize}

Selon le critère de Kaiser (\(\lambda > 1\)), les trois premières
dimensions sont retenues car leurs valeurs propres sont supérieures à 1.

\subsection{Graphiques de l'ACP}\label{graphiques-de-lacp}

\subsubsection{Diagramme des valeurs propres (Scree
Plot)}\label{diagramme-des-valeurs-propres-scree-plot}

\begin{Shaded}
\begin{Highlighting}[]
\FunctionTok{fviz\_eig}\NormalTok{(res.pca, }\AttributeTok{addlabels =} \ConstantTok{TRUE}\NormalTok{)}
\end{Highlighting}
\end{Shaded}

\begin{center}\includegraphics{rapport_ameliore_files/figure-latex/scree-plot-1} \end{center}

Le diagramme des valeurs propres montre que les premières composantes
principales concentrent une part importante de l'information. La
première dimension explique à elle seule \textbf{28,1 \%} de la variance
totale, et les trois premières composantes permettent de conserver plus
de \textbf{63 \%} de la variance, ce qui constitue une représentation
satisfaisante des données multidimensionnelles.

\subsubsection{Cercle des corrélations des
variables}\label{cercle-des-corruxe9lations-des-variables}

\begin{Shaded}
\begin{Highlighting}[]
\FunctionTok{fviz\_pca\_var}\NormalTok{(res.pca)}
\end{Highlighting}
\end{Shaded}

\begin{center}\includegraphics{rapport_ameliore_files/figure-latex/cercle-correlations-1} \end{center}

Le cercle des corrélations met en évidence plusieurs relations
significatives :

\begin{itemize}
\tightlist
\item
  \texttt{bmi} et \texttt{percent\_body\_fat} sont \textbf{fortement
  corrélées positivement}.
\item
  \texttt{vo2\_estimate\_ml\_per\_kg\_min} et \texttt{sit\_ups\_count}
  sont proches, traduisant une relation positive entre endurance
  cardiovasculaire et performance musculaire.
\item
  Une \textbf{opposition} est observée entre les variables de
  performance physique et les variables de composition corporelle.
\item
  La variable \texttt{age} présente une contribution plus modérée sur
  les deux premiers axes.
\end{itemize}

\subsubsection{Contribution des variables à la 1ère composante
principale}\label{contribution-des-variables-uxe0-la-1uxe8re-composante-principale}

\begin{Shaded}
\begin{Highlighting}[]
\FunctionTok{fviz\_contrib}\NormalTok{(res.pca, }\AttributeTok{choice =} \StringTok{"var"}\NormalTok{, }\AttributeTok{axes =} \DecValTok{1}\NormalTok{)}
\end{Highlighting}
\end{Shaded}

\begin{center}\includegraphics{rapport_ameliore_files/figure-latex/contrib-axe1-1} \end{center}

La première composante principale est principalement construite par les
variables liées à la condition physique générale :
\texttt{hand\_grip\_strength\_kg},
\texttt{vo2\_estimate\_ml\_per\_kg\_min} et \texttt{sit\_ups\_count}. La
variable \texttt{percent\_body\_fat} y apparaît opposée, indiquant qu'un
taux élevé de graisse corporelle est associé à de moins bonnes
performances physiques. Cette dimension peut être interprétée comme un
\textbf{axe global de condition physique et de santé}.

\subsubsection{Contribution des variables à la 2ème composante
principale}\label{contribution-des-variables-uxe0-la-2uxe8me-composante-principale}

\begin{Shaded}
\begin{Highlighting}[]
\FunctionTok{fviz\_contrib}\NormalTok{(res.pca, }\AttributeTok{choice =} \StringTok{"var"}\NormalTok{, }\AttributeTok{axes =} \DecValTok{2}\NormalTok{)}
\end{Highlighting}
\end{Shaded}

\begin{center}\includegraphics{rapport_ameliore_files/figure-latex/contrib-axe2-1} \end{center}

La deuxième composante principale est principalement liée à l'endurance
physique et aux caractéristiques physiologiques. Les variables
\texttt{age}, \texttt{vo2\_estimate\_ml\_per\_kg\_min} et
\texttt{sit\_ups\_count} présentent les contributions les plus
importantes. Cette dimension traduit une différenciation des individus
selon leur \textbf{capacité cardiovasculaire et leur profil
physiologique global}.

\subsubsection{Projection des individus sur le plan
factoriel}\label{projection-des-individus-sur-le-plan-factoriel}

\begin{Shaded}
\begin{Highlighting}[]
\FunctionTok{fviz\_pca\_ind}\NormalTok{(res.pca)}
\end{Highlighting}
\end{Shaded}

\begin{center}\includegraphics{rapport_ameliore_files/figure-latex/projection-individus-1} \end{center}

La projection des individus montre une dispersion relativement
importante des observations, traduisant une forte diversité des profils
physiques dans le dataset. Certains regroupements apparaissent
néanmoins, suggérant l'existence de profils homogènes, ce qui justifie
l'utilisation ultérieure du clustering K-Means.

\subsubsection{Qualité de représentation des variables
(Cos²)}\label{qualituxe9-de-repruxe9sentation-des-variables-cosuxb2}

\begin{Shaded}
\begin{Highlighting}[]
\FunctionTok{fviz\_cos2}\NormalTok{(res.pca, }\AttributeTok{choice =} \StringTok{"var"}\NormalTok{)}
\end{Highlighting}
\end{Shaded}

\begin{center}\includegraphics{rapport_ameliore_files/figure-latex/cos2-1} \end{center}

Les variables \texttt{hand\_grip\_strength\_kg},
\texttt{percent\_body\_fat} et \texttt{vo2\_estimate\_ml\_per\_kg\_min}
présentent les meilleurs cos², indiquant qu'elles sont correctement
représentées sur le plan factoriel. À l'inverse,
\texttt{sit\_and\_reach\_cm} présente une qualité de représentation plus
faible.

\subsubsection{Projection des individus selon le
sexe}\label{projection-des-individus-selon-le-sexe}

\begin{Shaded}
\begin{Highlighting}[]
\FunctionTok{fviz\_pca\_ind}\NormalTok{(}
\NormalTok{  res.pca,}
  \AttributeTok{habillage =}\NormalTok{ data}\SpecialCharTok{$}\NormalTok{sex}
\NormalTok{)}
\end{Highlighting}
\end{Shaded}

\begin{center}\includegraphics{rapport_ameliore_files/figure-latex/projection-sexe-1} \end{center}

L'ajout de la variable illustrative \texttt{sex} met en évidence des
différences partielles entre les profils masculins et féminins. La
séparation reste toutefois partielle, suggérant également des
caractéristiques communes entre les individus indépendamment du sexe.

\subsubsection{Matrice des
corrélations}\label{matrice-des-corruxe9lations}

\begin{Shaded}
\begin{Highlighting}[]
\FunctionTok{library}\NormalTok{(corrplot)}

\FunctionTok{corrplot}\NormalTok{(}
  \FunctionTok{cor}\NormalTok{(data\_scaled),}
  \AttributeTok{method    =} \StringTok{"color"}\NormalTok{,}
  \AttributeTok{type      =} \StringTok{"upper"}\NormalTok{,}
  \AttributeTok{addCoef.col =} \StringTok{"black"}\NormalTok{,}
  \AttributeTok{tl.col    =} \StringTok{"black"}\NormalTok{,}
  \AttributeTok{tl.cex    =} \FloatTok{0.7}\NormalTok{,}
  \AttributeTok{number.cex =} \FloatTok{0.6}
\NormalTok{)}
\end{Highlighting}
\end{Shaded}

\begin{center}\includegraphics{rapport_ameliore_files/figure-latex/matrice-correlations-1} \end{center}

La matrice des corrélations met en évidence :

\begin{itemize}
\tightlist
\item
  Une \textbf{forte corrélation positive} entre \texttt{bmi} et
  \texttt{percent\_body\_fat}.
\item
  Des \textbf{corrélations positives} entre les variables de performance
  physique (\texttt{vo2\_estimate\_ml\_per\_kg\_min},
  \texttt{sit\_ups\_count}, \texttt{hand\_grip\_strength\_kg}).
\item
  Des \textbf{corrélations négatives} entre composition corporelle et
  performance physique.
\end{itemize}

Cette structure justifie pleinement l'utilisation de l'ACP pour résumer
les relations multidimensionnelles entre les variables.

\subsection{Conclusion de l'ACP}\label{conclusion-de-lacp}

L'Analyse en Composantes Principales a permis de mettre en évidence les
principales structures statistiques des données. Les résultats montrent
que :

\begin{itemize}
\tightlist
\item
  La \textbf{1ère composante} représente essentiellement le niveau
  global de condition physique.
\item
  La \textbf{2ème composante} est principalement liée à l'âge.
\item
  La \textbf{masse grasse} est négativement associée aux performances
  physiques.
\item
  Les \textbf{capacités musculaires et cardiovasculaires} sont
  positivement corrélées.
\item
  Des \textbf{différences entre les sexes} existent mais restent
  partiellement superposées.
\end{itemize}

\begin{center}\rule{0.5\linewidth}{0.5pt}\end{center}

\section{Classification des Individus --- Clustering
K-Means}\label{classification-des-individus-clustering-k-means}

\subsection{Objectif du clustering}\label{objectif-du-clustering}

Après la réduction de dimension par l'ACP, une méthode de classification
non supervisée a été appliquée afin d'identifier des groupes homogènes
d'individus partageant des caractéristiques physiques similaires.
L'algorithme \textbf{K-Means} a été retenu pour sa capacité à
partitionner efficacement les individus en maximisant la similarité
intra-cluster et la dissimilarité inter-clusters.

\begin{Shaded}
\begin{Highlighting}[]
\FunctionTok{library}\NormalTok{(cluster)}
\FunctionTok{library}\NormalTok{(factoextra)}
\end{Highlighting}
\end{Shaded}

\subsection{Choix du nombre de clusters --- Méthode du
coude}\label{choix-du-nombre-de-clusters-muxe9thode-du-coude}

Avant l'application du K-Means, la \textbf{méthode du coude}
(\emph{Elbow Method}) a été utilisée pour déterminer le nombre optimal
de clusters, en analysant l'évolution de l'inertie intra-cluster en
fonction de K.

\begin{Shaded}
\begin{Highlighting}[]
\FunctionTok{fviz\_nbclust}\NormalTok{(}
\NormalTok{  data\_scaled,}
\NormalTok{  kmeans,}
  \AttributeTok{method =} \StringTok{"wss"}
\NormalTok{) }\SpecialCharTok{+}
  \FunctionTok{labs}\NormalTok{(}\AttributeTok{title =} \StringTok{"Méthode du coude pour le choix du nombre de clusters"}\NormalTok{)}
\end{Highlighting}
\end{Shaded}

\begin{center}\includegraphics{rapport_ameliore_files/figure-latex/elbow-1} \end{center}

L'analyse du graphique montre une diminution importante de l'inertie
jusqu'à \textbf{K = 3}, après quoi la réduction devient plus faible. Le
choix de \textbf{trois clusters} apparaît donc comme un compromis
satisfaisant entre simplicité et qualité de segmentation.

\subsection{Application de K-Means}\label{application-de-k-means}

L'algorithme K-Means est appliqué aux données standardisées avec trois
groupes et le paramètre \texttt{nstart\ =\ 25} pour améliorer la
stabilité de la solution.

\begin{Shaded}
\begin{Highlighting}[]
\FunctionTok{set.seed}\NormalTok{(}\DecValTok{123}\NormalTok{)}

\NormalTok{kmeans\_result }\OtherTok{\textless{}{-}} \FunctionTok{kmeans}\NormalTok{(}
\NormalTok{  data\_scaled,}
  \AttributeTok{centers =} \DecValTok{3}\NormalTok{,}
  \AttributeTok{nstart  =} \DecValTok{25}
\NormalTok{)}

\NormalTok{kmeans\_result}
\end{Highlighting}
\end{Shaded}

\begin{verbatim}
## K-means clustering with 3 clusters of sizes 617, 712, 671
## 
## Cluster means:
##           age        bmi percent_body_fat hand_grip_strength_kg
## 1  0.08329186 -0.2055530        1.2954349            -1.0070737
## 2  0.61296942  0.4663638       -0.4985752             0.6944622
## 3 -0.72701238 -0.3058493       -0.6621428             0.1891317
##   sit_and_reach_cm sit_ups_count vo2_estimate_ml_per_kg_min
## 1       0.22711787     0.3967906                -0.69566896
## 2      -0.25049174    -0.4538260                -0.09412882
## 3       0.05695736     0.1166979                 0.73956403
## 
## Clustering vector:
##    [1] 2 3 2 3 1 1 1 3 2 2 1 2 3 1 1 1 1 2 2 1 2 1 1 1 3 2 1 3 2 1 3 1 1 2 3 3 1
##   [38] 1 2 2 1 2 1 3 1 3 1 2 2 1 2 3 2 2 1 3 1 1 3 2 3 1 2 2 1 2 1 2 1 2 3 1 1 2
##   [75] 3 3 3 1 3 1 3 3 2 1 1 3 3 2 2 2 1 2 2 3 2 2 3 3 1 1 1 2 1 2 2 1 2 3 1 1 1
##  [112] 1 3 2 3 3 2 1 2 2 3 2 1 1 1 3 3 2 1 3 3 1 1 2 2 2 2 3 3 2 2 1 2 1 1 1 3 3
##  [149] 3 1 3 1 3 2 2 1 2 2 1 3 3 2 3 2 1 3 2 1 1 3 3 1 3 1 2 1 3 2 3 1 3 1 3 3 1
##  [186] 2 3 3 3 1 1 3 2 3 2 2 3 2 3 2 3 1 1 3 2 1 1 2 1 1 3 2 2 1 3 1 2 2 2 3 2 2
##  [223] 1 3 1 1 2 3 3 3 2 2 3 2 1 2 1 1 2 1 3 3 2 3 1 2 2 2 3 3 1 3 3 1 1 3 3 2 3
##  [260] 1 2 2 1 2 3 3 2 3 2 1 3 3 3 3 3 2 3 2 2 3 3 2 3 1 2 1 3 3 1 2 1 1 2 2 1 2
##  [297] 3 1 3 1 1 3 3 3 2 3 2 3 3 1 2 1 1 2 3 3 2 1 3 1 2 2 2 1 2 3 3 2 3 1 2 2 3
##  [334] 1 1 1 2 3 2 1 2 2 1 3 2 2 3 1 2 3 3 2 2 1 1 3 1 3 3 2 3 2 3 3 2 3 2 3 1 1
##  [371] 1 1 1 2 1 1 3 1 2 1 1 3 2 2 3 2 1 1 2 2 3 2 3 3 2 2 3 2 2 2 1 2 3 2 2 3 3
##  [408] 1 3 1 1 3 3 3 2 2 1 3 2 2 2 2 1 2 1 2 2 2 3 1 1 2 2 1 1 2 3 2 2 1 1 3 2 3
##  [445] 3 3 2 2 3 1 2 3 2 3 1 1 1 2 1 1 1 2 3 3 3 2 1 2 1 2 1 2 1 3 3 3 1 3 3 2 2
##  [482] 2 2 2 3 2 1 3 2 2 3 2 1 1 3 3 3 1 2 3 3 2 1 3 3 1 2 3 2 2 3 2 2 2 1 1 2 2
##  [519] 3 1 3 2 3 3 2 2 2 1 1 2 2 3 3 2 3 2 2 1 3 1 2 2 1 2 3 3 1 2 3 2 3 3 3 2 2
##  [556] 2 1 1 1 3 2 1 1 1 3 3 3 2 2 3 3 2 2 2 2 2 1 2 3 1 2 3 2 3 1 1 1 3 1 1 1 2
##  [593] 1 3 2 2 2 1 1 1 1 1 1 1 1 3 1 2 2 3 1 3 3 1 2 3 1 1 2 3 1 1 2 2 3 2 2 3 1
##  [630] 1 2 2 2 2 1 2 2 1 1 3 1 3 1 2 1 2 3 2 3 2 1 2 3 2 2 1 3 2 2 3 3 2 3 1 1 1
##  [667] 1 3 1 1 1 2 1 3 2 2 1 3 3 3 2 1 1 2 2 2 2 1 2 1 2 2 3 1 1 3 1 3 1 2 2 1 3
##  [704] 1 1 3 2 3 3 3 2 2 3 2 2 3 2 2 1 3 3 1 1 2 2 3 1 2 1 1 1 2 2 2 3 3 2 2 2 3
##  [741] 3 1 3 3 1 2 3 2 3 3 1 2 3 2 1 3 2 3 2 3 1 1 1 2 1 3 2 3 3 2 2 2 1 2 2 1 1
##  [778] 2 2 3 2 3 1 3 1 2 3 2 2 3 1 3 2 2 1 2 3 3 3 3 3 1 2 3 3 2 3 3 2 1 2 3 1 2
##  [815] 2 3 3 3 1 2 3 1 2 3 3 2 3 3 3 2 1 3 1 1 2 3 2 2 3 3 1 3 1 2 1 3 1 3 2 3 1
##  [852] 3 1 3 1 1 1 2 2 1 3 3 2 3 2 2 1 1 2 2 3 1 2 3 1 3 2 2 3 1 1 3 1 1 2 1 2 3
##  [889] 1 3 2 2 2 2 2 1 1 2 1 1 1 1 1 2 2 3 3 3 3 3 3 1 2 3 3 2 1 3 2 1 2 2 3 2 3
##  [926] 1 1 1 3 1 1 3 2 3 1 1 1 3 1 3 1 3 3 2 3 1 1 2 2 2 2 2 1 3 2 3 3 1 1 2 1 2
##  [963] 1 2 1 2 2 1 2 3 3 2 3 1 3 1 1 1 3 3 2 2 1 3 2 2 3 2 1 1 2 2 1 2 2 1 2 1 3
## [1000] 3 1 1 2 2 3 2 2 2 1 2 1 3 3 1 2 3 1 3 2 1 2 1 2 2 3 3 3 2 3 2 1 3 1 1 1 2
## [1037] 3 3 1 1 2 3 1 2 2 3 1 1 2 1 3 3 2 3 3 2 3 3 2 2 2 3 2 3 3 2 1 1 3 1 1 3 1
## [1074] 1 2 3 2 3 1 2 3 1 3 2 3 3 2 1 2 1 2 2 3 2 1 1 1 1 3 3 1 3 3 3 3 1 3 3 2 1
## [1111] 2 1 1 3 3 2 3 2 1 1 2 2 3 3 3 3 1 2 2 1 2 3 3 1 2 1 3 2 2 2 3 2 2 3 1 3 3
## [1148] 2 2 3 3 2 3 2 3 2 3 3 2 1 3 1 2 3 1 2 2 3 1 1 2 2 2 1 3 3 3 1 2 2 3 2 2 2
## [1185] 1 2 2 1 1 2 1 2 2 2 3 2 3 2 3 1 3 1 1 3 2 3 2 3 1 3 2 2 1 3 2 1 1 3 2 1 2
## [1222] 1 2 2 3 3 3 3 2 2 3 1 1 1 2 2 3 3 2 1 2 2 2 2 3 3 1 1 2 3 3 3 1 1 3 3 1 3
## [1259] 1 3 1 2 1 2 1 1 2 3 2 1 2 2 3 1 2 3 1 2 1 3 3 3 2 1 2 3 1 3 3 1 3 1 2 3 3
## [1296] 1 1 1 2 3 1 3 1 3 2 2 3 2 3 2 1 2 3 2 2 2 3 3 1 1 1 1 2 1 2 2 1 2 3 2 3 2
## [1333] 1 1 3 3 1 3 1 3 1 3 3 3 3 2 2 1 2 1 2 2 3 2 2 3 2 3 1 2 3 1 2 2 3 1 2 2 1
## [1370] 1 3 1 2 3 1 3 2 3 3 3 1 2 2 2 2 3 3 1 2 2 1 1 1 2 2 3 2 1 2 3 1 1 1 3 2 2
## [1407] 3 3 2 2 1 1 1 1 2 3 1 1 3 3 1 3 1 1 2 3 1 3 3 2 3 2 3 3 1 1 1 3 1 3 3 1 2
## [1444] 3 2 1 3 3 2 2 2 3 3 2 1 3 2 3 3 3 3 3 1 2 2 1 3 2 2 3 3 2 3 2 3 3 3 3 3 3
## [1481] 2 1 3 1 2 3 2 2 2 3 1 3 1 1 2 3 3 2 3 2 3 2 1 2 3 3 2 2 2 3 3 2 3 1 1 1 2
## [1518] 2 1 3 1 3 3 1 3 1 1 1 2 2 1 3 3 1 3 3 1 3 1 3 2 1 3 1 2 2 2 2 3 2 2 3 3 2
## [1555] 1 1 3 2 3 2 2 1 2 1 3 2 3 1 2 3 3 1 3 3 1 1 1 2 2 3 1 2 2 1 2 1 2 2 2 2 2
## [1592] 2 1 1 2 2 1 1 3 2 1 3 1 3 1 2 1 3 1 3 2 3 3 3 3 2 2 3 2 2 2 1 3 2 1 2 1 2
## [1629] 1 3 2 3 3 2 2 1 3 1 2 3 3 2 2 3 1 2 1 2 3 2 1 3 1 3 3 3 3 1 3 2 2 2 2 3 3
## [1666] 2 1 3 1 2 3 3 2 3 3 1 1 3 1 2 2 2 3 2 1 2 2 2 2 3 1 2 2 1 2 1 1 1 3 3 2 2
## [1703] 1 2 2 1 1 1 2 3 1 1 3 2 1 3 1 2 2 1 1 1 1 1 2 1 3 3 2 3 1 3 2 3 1 2 2 1 2
## [1740] 2 3 3 2 1 3 2 1 1 3 1 3 1 1 1 2 2 2 2 1 3 3 3 1 2 1 2 2 3 3 2 2 3 2 2 2 1
## [1777] 2 2 1 3 1 3 3 2 3 2 2 3 1 3 1 2 2 3 1 1 2 1 1 1 3 3 2 1 2 3 3 1 1 1 3 2 1
## [1814] 1 1 3 1 2 2 2 3 1 3 2 3 3 1 2 3 1 2 3 3 3 2 3 1 2 2 1 3 2 3 2 2 3 1 1 2 3
## [1851] 1 2 2 3 3 1 2 3 3 2 2 3 1 2 2 2 1 1 1 2 3 1 2 3 2 3 2 3 1 2 3 3 2 3 1 2 1
## [1888] 1 2 2 1 3 2 3 3 2 2 1 2 2 3 1 2 3 3 1 3 3 3 1 1 3 3 3 2 2 2 2 3 3 3 1 3 2
## [1925] 2 1 1 2 2 1 2 2 3 3 3 2 1 3 3 3 1 3 2 3 2 2 3 1 3 3 1 1 3 3 3 2 2 2 3 2 1
## [1962] 2 2 2 3 3 1 3 2 1 1 3 1 3 1 2 1 3 2 2 3 2 3 3 1 1 1 2 3 3 3 1 3 1 3 2 2 1
## [1999] 1 1
## 
## Within cluster sum of squares by cluster:
## [1] 3203.552 3148.571 3267.472
##  (between_SS / total_SS =  31.3 %)
## 
## Available components:
## 
## [1] "cluster"      "centers"      "totss"        "withinss"     "tot.withinss"
## [6] "betweenss"    "size"         "iter"         "ifault"
\end{verbatim}

\begin{Shaded}
\begin{Highlighting}[]
\NormalTok{health\_clean}\SpecialCharTok{$}\NormalTok{cluster }\OtherTok{\textless{}{-}} \FunctionTok{as.factor}\NormalTok{(kmeans\_result}\SpecialCharTok{$}\NormalTok{cluster)}
\end{Highlighting}
\end{Shaded}

L'algorithme K-Means a segmenté les 2 000 individus en trois groupes aux
tailles relativement équilibrées :

{\def\LTcaptype{none} % do not increment counter
\begin{longtable}[]{@{}ll@{}}
\toprule\noalign{}
Cluster & Nombre d'individus \\
\midrule\noalign{}
\endhead
\bottomrule\noalign{}
\endlastfoot
Cluster 1 & 617 \\
Cluster 2 & 712 \\
Cluster 3 & 671 \\
\end{longtable}
}

\subsection{Visualisation des
clusters}\label{visualisation-des-clusters}

\begin{Shaded}
\begin{Highlighting}[]
\FunctionTok{fviz\_cluster}\NormalTok{(}
\NormalTok{  kmeans\_result,}
  \AttributeTok{data         =}\NormalTok{ data\_scaled,}
  \AttributeTok{ellipse.type =} \StringTok{"norm"}\NormalTok{,}
  \AttributeTok{geom         =} \StringTok{"point"}\NormalTok{,}
  \AttributeTok{palette      =} \StringTok{"jco"}\NormalTok{,}
  \AttributeTok{ggtheme      =} \FunctionTok{theme\_minimal}\NormalTok{()}
\NormalTok{)}
\end{Highlighting}
\end{Shaded}

\begin{center}\includegraphics{rapport_ameliore_files/figure-latex/visu-clusters-1} \end{center}

Les résultats montrent une structuration relativement claire des
clusters, confirmant la pertinence de la classification obtenue.

\subsection{Analyse des profils par
cluster}\label{analyse-des-profils-par-cluster}

Les moyennes des variables quantitatives sont calculées pour chaque
cluster afin d'interpréter les caractéristiques physiques et
physiologiques de chaque groupe.

\begin{Shaded}
\begin{Highlighting}[]
\NormalTok{profil\_clusters }\OtherTok{\textless{}{-}} \FunctionTok{aggregate}\NormalTok{(}
\NormalTok{  health\_clean[, variables\_quantitatives],}
  \AttributeTok{by  =} \FunctionTok{list}\NormalTok{(}\AttributeTok{Cluster =}\NormalTok{ health\_clean}\SpecialCharTok{$}\NormalTok{cluster),}
  \AttributeTok{FUN =}\NormalTok{ mean}
\NormalTok{)}

\NormalTok{profil\_clusters}
\end{Highlighting}
\end{Shaded}

\begin{verbatim}
##   Cluster      age      bmi percent_body_fat hand_grip_strength_kg
## 1       1 42.64182 23.31831         56.80810              25.14489
## 2       2 49.63904 25.47626         44.35304              43.41297
## 3       3 31.93741 22.99620         43.21746              37.98763
##   sit_and_reach_cm sit_ups_count vo2_estimate_ml_per_kg_min
## 1         9.264344      19.36305                   31.37423
## 2         6.367697      15.62640                   35.61868
## 3         8.232340      18.13264                   41.50119
\end{verbatim}

\textbf{Cluster 1 --- Profil à faible condition physique}

Ce groupe se caractérise principalement par un pourcentage de graisse
corporelle très élevé, une faible force musculaire et des capacités
cardiovasculaires réduites (VO₂ max le plus faible parmi les trois
groupes). Ce profil peut correspondre à des individus physiquement moins
actifs ou présentant un risque accru lié à la sédentarité.

\textbf{Cluster 2 --- Profil musculaire intermédiaire}

Ce deuxième groupe correspond à des individus relativement plus âgés
présentant une bonne force musculaire (force de préhension la plus
élevée) et un indice de masse corporelle légèrement supérieur à la
moyenne. Malgré un âge moyen plus élevé, les individus conservent des
performances musculaires satisfaisantes, mais avec une endurance
cardiovasculaire plus modérée.

\textbf{Cluster 3 --- Profil sportif et performant}

Ce troisième cluster regroupe les individus présentant la meilleure
condition physique globale : âge moyen plus jeune, faible IMC, faible
taux de graisse corporelle et meilleures capacités cardiovasculaires.
Ces individus présentent globalement un profil plus sportif et
physiquement performant.

\subsection{Distribution des variables par
cluster}\label{distribution-des-variables-par-cluster}

\begin{Shaded}
\begin{Highlighting}[]
\FunctionTok{ggplot}\NormalTok{(health\_clean, }\FunctionTok{aes}\NormalTok{(}\AttributeTok{x =}\NormalTok{ cluster, }\AttributeTok{y =}\NormalTok{ bmi, }\AttributeTok{fill =}\NormalTok{ cluster)) }\SpecialCharTok{+}
  \FunctionTok{geom\_boxplot}\NormalTok{() }\SpecialCharTok{+}
  \FunctionTok{theme\_minimal}\NormalTok{() }\SpecialCharTok{+}
  \FunctionTok{labs}\NormalTok{(}
    \AttributeTok{title =} \StringTok{"Distribution du BMI par cluster"}\NormalTok{,}
    \AttributeTok{x     =} \StringTok{"Cluster"}\NormalTok{,}
    \AttributeTok{y     =} \StringTok{"IMC"}
\NormalTok{  )}
\end{Highlighting}
\end{Shaded}

\begin{center}\includegraphics{rapport_ameliore_files/figure-latex/boxplots-clusters-1} \end{center}

\begin{Shaded}
\begin{Highlighting}[]
\FunctionTok{ggplot}\NormalTok{(health\_clean, }\FunctionTok{aes}\NormalTok{(}\AttributeTok{x =}\NormalTok{ cluster, }\AttributeTok{y =}\NormalTok{ vo2\_estimate\_ml\_per\_kg\_min, }\AttributeTok{fill =}\NormalTok{ cluster)) }\SpecialCharTok{+}
  \FunctionTok{geom\_boxplot}\NormalTok{() }\SpecialCharTok{+}
  \FunctionTok{theme\_minimal}\NormalTok{() }\SpecialCharTok{+}
  \FunctionTok{labs}\NormalTok{(}
    \AttributeTok{title =} \StringTok{"Distribution du VO₂ max par cluster"}\NormalTok{,}
    \AttributeTok{x     =} \StringTok{"Cluster"}\NormalTok{,}
    \AttributeTok{y     =} \StringTok{"VO₂ max (ml/kg/min)"}
\NormalTok{  )}
\end{Highlighting}
\end{Shaded}

\begin{center}\includegraphics{rapport_ameliore_files/figure-latex/boxplots-clusters-2} \end{center}

Les boxplots confirment des différences significatives entre les
clusters en termes de composition corporelle et de performance
cardiovasculaire, validant la capacité du modèle K-Means à segmenter
efficacement les individus.

\subsection{Conclusion du clustering}\label{conclusion-du-clustering}

L'algorithme K-Means a permis d'identifier trois profils distincts
d'individus à partir des variables biométriques et physiques étudiées.
Les résultats montrent que les données présentent une structure
naturelle permettant de distinguer différents niveaux de condition
physique au sein de la population étudiée. Les clusters obtenus pourront
servir de base pour des analyses décisionnelles ou des stratégies de
suivi adaptées aux différents profils identifiés.

\begin{center}\rule{0.5\linewidth}{0.5pt}\end{center}

\section{Analyse Combinée --- ACP et
Clustering}\label{analyse-combinuxe9e-acp-et-clustering}

\subsection{Projection des clusters sur le plan
factoriel}\label{projection-des-clusters-sur-le-plan-factoriel}

Les clusters obtenus par K-Means sont projetés sur le plan factoriel
défini par les deux premières composantes principales de l'ACP,
permettant d'observer leur répartition dans l'espace réduit.

\begin{Shaded}
\begin{Highlighting}[]
\FunctionTok{fviz\_cluster}\NormalTok{(}
\NormalTok{  kmeans\_result,}
  \AttributeTok{data         =}\NormalTok{ data\_scaled,}
  \AttributeTok{ellipse.type =} \StringTok{"norm"}\NormalTok{,}
  \AttributeTok{geom         =} \StringTok{"point"}\NormalTok{,}
  \AttributeTok{palette      =} \StringTok{"jco"}\NormalTok{,}
  \AttributeTok{ggtheme      =} \FunctionTok{theme\_minimal}\NormalTok{()}
\NormalTok{)}
\end{Highlighting}
\end{Shaded}

\begin{center}\includegraphics{rapport_ameliore_files/figure-latex/projection-clusters-1} \end{center}

La projection montre une séparation relativement claire entre les
groupes, confirmant que les dimensions extraites par l'ACP capturent
efficacement les principales différences entre les profils physiques.
Quelques zones de chevauchement restent observables, ce qui est fréquent
dans des données multidimensionnelles réelles.

\subsection{Relation entre ACP et
segmentation}\label{relation-entre-acp-et-segmentation}

L'ACP et le clustering K-Means sont deux méthodes complémentaires. L'ACP
réduit la dimensionnalité des données tout en conservant l'essentiel de
l'information statistique, facilitant ainsi la segmentation des
individus en éliminant une partie du bruit et des redondances entre
variables. Les axes factoriels représentent les principales sources de
variabilité utilisées indirectement par l'algorithme de clustering.

\subsection{Importance des variables dans la
segmentation}\label{importance-des-variables-dans-la-segmentation}

Les variables jouant un rôle particulièrement important dans la
différenciation des clusters sont :

\begin{itemize}
\tightlist
\item
  \textbf{Variables de performance physique} :
  \texttt{vo2\_estimate\_ml\_per\_kg\_min},
  \texttt{hand\_grip\_strength\_kg}, \texttt{sit\_ups\_count}
\item
  \textbf{Variables de composition corporelle} : \texttt{bmi},
  \texttt{percent\_body\_fat} (permettent d'identifier les profils à
  condition physique plus faible)
\end{itemize}

Ces résultats confirment les observations réalisées lors de l'ACP.

\subsection{Conclusion de l'analyse
combinée}\label{conclusion-de-lanalyse-combinuxe9e}

L'analyse combinée de l'ACP et du clustering K-Means a permis d'obtenir
une vision globale et structurée des données. L'ACP a facilité la
compréhension des relations entre les variables et réduit la complexité
du dataset, tandis que le clustering a permis d'identifier plusieurs
profils homogènes d'individus.

Les résultats montrent une \textbf{cohérence importante} entre les axes
factoriels et les groupes identifiés. Cette approche combinée permet
ainsi de mieux comprendre la structure multidimensionnelle des données
et de révéler différents niveaux de condition physique au sein de la
population étudiée.

\begin{center}\rule{0.5\linewidth}{0.5pt}\end{center}

\emph{Fin du rapport --- Analyse réalisée avec R.}

\end{document}
